\chapter{Représentations des groupes finis}\label{ch:b3:representations}

Pour comprendre un groupe abstrait, le faire agir sur un espace
vectoriel et utiliser l'algèbre linéaire --- valeurs propres,
traces, produits scalaires --- sur l'\omterm{def:b3:groups:action}{action}. Ce programme, la
\emph{théorie des \omterm{def:b3:representations:rep}{représentations}}, est étonnamment efficace pour
les groupes finis sur $\C$~: toute \omterm{def:b3:representations:rep}{représentation} se scinde en
\omterm{def:b3:representations:rep}{irréductibles} (Maschke), les \omterm{def:b3:representations:rep}{irréductibles} sont repérées par leurs
\emph{\omterm{def:b3:representations:character}{caractères}} (traces), et les \omterm{def:b3:representations:character}{caractères} satisfont des
relations d'orthogonalité qui rendent les calculs mécaniques. Le
chapitre construit ce calcul et ses premiers chefs-d'œuvre ---
\omterm{def:b3:representations:table}{tables de caractères} des petits groupes --- et le problème du
week-end récolte un théorème bien au-delà de la pure théorie des
groupes~: le théorème $p^aq^b$ de Burnside. Tout au long, $G$ est
un groupe fini et tous les espaces vectoriels sont de dimension
finie sur $\C$.

\section{Représentations, Maschke, Schur}

\begin{definition}\label{def:b3:representations:rep}
Une \emph{représentation}\index{représentation} de $G$ est un
morphisme $\rho \colon G \to GL(V)$ pour un $\C$-espace vectoriel
$V$~; $\dim V$ est son \emph{\omterm{def:b3:galois:extension}{degré}}. Un sous-espace $W \subseteq
V$ est \emph{invariant} si $\rho(g)W \subseteq W$ pour tout $g$~;
la restriction fait de $W$ une \emph{sous-représentation}. $\rho$
est \emph{irréductible}\index{représentation irréductible} si $V
\neq 0$ et ses seuls sous-espaces invariants sont $0$ et $V$. Un
\emph{morphisme} entre $(\rho, V)$ et $(\sigma, W)$ est une
application linéaire $f\colon V \to W$ avec $f\rho(g) =
\sigma(g)f$ pour tout $g$ (équivariance)~; les $f$ bijectifs sont
des \emph{isomorphismes}.
\end{definition}

\begin{example}\label{ex:b3:representations:examples}
(a) \omterm{def:b3:galois:extension}{Degré} $1$~: morphismes $G \to \C^\times$.
(b) La \emph{représentation régulière}\index{représentation
régulière}~: $V = \C^G$ de base $(e_h)_{h \in G}$,
$\rho(g)e_h = e_{gh}$~; \omterm{def:b3:galois:extension}{degré} $\abs G$.
(c) Une \omterm{def:b3:groups:action}{action} de permutation de $G$ sur un ensemble fini $X$
donne la \emph{\omterm{def:b3:representations:rep}{représentation} de permutation} sur $\C^X$~:
$\rho(g)e_x = e_{g\cdot x}$.
(d) $S_n$ agit sur $\C^n$ en permutant les coordonnées~;
l'hyperplan $\{\sum x_i = 0\}$ est invariant~: la
\emph{\omterm{def:b3:representations:rep}{représentation} standard}, de \omterm{def:b3:galois:extension}{degré} $n - 1$.
\end{example}

\begin{theorem}[Maschke]\label{thm:b3:representations:maschke}
Tout sous-espace invariant $W$ d'une \omterm{def:b3:representations:rep}{représentation} $(V, \rho)$
admet un supplémentaire invariant. Par conséquent toute
\omterm{def:b3:representations:rep}{représentation} est somme directe d'\omterm{def:b3:representations:rep}{irréductibles}
(\emph{semisimplicité}).\index{théorème de Maschke}
\end{theorem}

\begin{proof}
Soit $p \colon V \to V$ \emph{une} projection d'image $W$
(choisir un quelconque supplémentaire). La moyenner sur le
groupe~:
\[
\tilde p = \frac1{\abs G}\sum_{g \in G} \rho(g)\,p\,\rho(g)^{-1}.
\]
Chaque terme envoie $V$ dans $W$ ($W$ invariant), et fixe $W$
point par point~: pour $w \in W$, $\rho(g)^{-1}w \in W$, $p$ le
fixe, et $\rho(g)$ le ramène --- donc $\tilde p$ est encore une
projection sur $W$. Elle est équivariante~: pour $h \in G$,
$\rho(h)\tilde p\rho(h)^{-1}$ réindexe la même somme. Donc $\ker
\tilde p$ est un supplémentaire invariant de $W$. En itérant sur
les facteurs (dimension finie) on décompose $V$ en \omterm{def:b3:representations:rep}{irréductibles}.
\end{proof}

\begin{theorem}[Lemme de Schur]\label{thm:b3:representations:schur}
Soit $f \colon V \to W$ un morphisme de \omterm{def:b3:representations:rep}{représentations}
\emph{\omterm{def:b3:representations:rep}{irréductibles}}. Alors $f = 0$ ou $f$ est un isomorphisme~;
et si $(V,\rho) = (W,\sigma)$, alors $f = \lambda\,\mathrm{id}$
pour un $\lambda \in \C$. D'où $\dim\operatorname{Hom}_G(V, W)$
vaut $1$ si $V \cong W$, $0$ sinon.\index{lemme de Schur}
\end{theorem}

\begin{proof}
$\ker f$ et $\operatorname{im} f$ sont invariants
(équivariance), donc chacun est $0$ ou tout~: soit $f = 0$, soit
$f$ est injectif d'image pleine. Si $V = W$~: $f$ a une valeur
propre $\lambda$ ($\C$ \omterm{def:b3:galois:closure}{algébriquement clos})~; $f -
\lambda\,\mathrm{id}$ est un morphisme non injectif $V \to V$,
donc $0$. Pour le dénombrement de dimension quand $V \cong W$~:
en fixant un isomorphisme $u$, tout morphisme $f$ donne
l'endomorphisme $u^{-1}f = \lambda\,\mathrm{id}$~: $f =
\lambda u$.
\end{proof}

\section{Caractères et orthogonalité}

\begin{definition}\label{def:b3:representations:character}
Le \emph{caractère}\index{caractère} de $(V, \rho)$ est
$\chi_\rho(g) = \operatorname{tr}\rho(g)$. Il vérifie
$\chi_\rho(e) = \dim V$, $\chi_\rho(hgh^{-1}) = \chi_\rho(g)$
(les traces sont invariantes par conjugaison)~: les caractères
sont des \emph{fonctions de classe}\index{fonction de classe}
--- éléments de l'espace $\mathcal{CF}(G)$ des fonctions
constantes sur les \omterm{ex:b3:groups:actions}{classes de conjugaison}, muni du produit
scalaire hermitien
\[
\langle \varphi, \psi\rangle = \frac1{\abs G}\sum_{g \in G}
\overline{\varphi(g)}\,\psi(g).
\]
\end{definition}

\begin{proposition}\label{prop:b3:representations:charbasics}
$\rho(g)$ est diagonalisable à valeurs propres racines de
l'unité~; $\chi_\rho(g^{-1}) = \overline{\chi_\rho(g)}$, et
$\abs{\chi_\rho(g)} \leq \chi_\rho(e)$ avec égalité ssi
$\rho(g)$ est scalaire. Les \omterm{def:b3:representations:character}{caractères} s'ajoutent sur les
sommes directes~: $\chi_{V \oplus W} = \chi_V + \chi_W$.
\end{proposition}

\begin{proof}
$\rho(g)^N = \mathrm{id}$ pour $N = \abs G$ (Lagrange)~:
$\rho(g)$ annule $X^N - 1$, scindé à racines \omterm{def:b3:groups:simple}{simples}, donc
diagonalisable à valeurs propres $\lambda_i \in \mu_N$. Alors
$\chi(g^{-1}) = \sum \lambda_i^{-1} = \sum\bar\lambda_i =
\overline{\chi(g)}$, et $\abs{\chi(g)} = \abs{\sum\lambda_i}
\leq \dim V$, avec égalité dans l'inégalité triangulaire ssi
toutes les $\lambda_i$ sont égales, i.e.\ $\rho(g) =
\lambda\,\mathrm{id}$. Additivité sur les sommes directes~:
traces par blocs.
\end{proof}

\begin{lemma}\label{lem:b3:representations:homtrace}
Soient $(V, \rho)$, $(W, \sigma)$ des \omterm{def:b3:representations:rep}{représentations}.
L'opérateur de moyenne sur $\operatorname{Hom}(W, V)$,
\[
c(A) = \frac1{\abs G}\sum_{g}\rho(g)\,A\,\sigma(g)^{-1},
\]
est une projection sur $\operatorname{Hom}_G(W, V)$, et pour
l'application $\Phi_g \colon A \mapsto \rho(g)A\sigma(g)^{-1}$
on a $\operatorname{tr}\Phi_g =
\chi_\rho(g)\,\overline{\chi_\sigma(g)}$.
\end{lemma}

\begin{proof}
$c(A)$ est équivariant (réindexer la somme comme chez Maschke),
et $c$ fixe les applications équivariantes (chaque terme égale
$A$)~: $c$ est une projection d'image $\operatorname{Hom}_G(W,
V)$. Trace~: en bases, $\Phi_g(A) = BAC$ avec $B = \rho(g)$, $C
= \sigma(g)^{-1}$~; sur la base $(E_{kl})$ des matrices,
$BE_{kl}C = \sum_{m,n} b_{mk}c_{ln}E_{mn}$, donc le coefficient
de $E_{kl}$ dans $\Phi_g(E_{kl})$ est $b_{kk}c_{ll}$~:
$\operatorname{tr}\Phi_g = \sum_{k,l}b_{kk}c_{ll} =
\operatorname{tr}(B) \operatorname{tr}(C) =
\chi_\rho(g)\chi_\sigma(g^{-1})$, et $\chi_\sigma(g^{-1}) =
\overline{\chi_\sigma(g)}$.
\end{proof}

\begin{theorem}[Premières relations d'orthogonalité]
\label{thm:b3:representations:orthogonality}
Soient $\rho, \sigma$ \emph{\omterm{def:b3:representations:rep}{irréductibles}}. Alors\index{orthogonalité
des caractères}
\[
\langle\chi_\sigma, \chi_\rho\rangle =
\begin{cases}
1 & \text{si } \rho \cong \sigma,\\
0 & \text{sinon~:}
\end{cases}
\]
les \omterm{def:b3:representations:character}{caractères} \omterm{def:b3:representations:rep}{irréductibles} forment une famille orthonormée
dans $\mathcal{CF}(G)$.
\end{theorem}

\begin{proof}
La trace d'une projection est la dimension de son image~:
\[
\dim\operatorname{Hom}_G(W, V) = \operatorname{tr} c
= \frac1{\abs G}\sum_g \operatorname{tr}\Phi_g
= \frac1{\abs G}\sum_g \chi_\rho(g)\overline{\chi_\sigma(g)}
= \langle \chi_\sigma, \chi_\rho\rangle,
\]
et le lemme de Schur évalue le premier membre à $\delta_{\rho
\cong \sigma}$.
\end{proof}

\begin{corollary}\label{cor:b3:representations:multiplicity}
Décomposer $V \cong \bigoplus_i V_i^{\oplus m_i}$ en
\omterm{def:b3:representations:rep}{irréductibles} distincts ($V_i \not\cong V_j$). Alors $m_i =
\langle \chi_{V_i}, \chi_V\rangle$~: les multiplicités --- donc
la \omterm{def:b3:representations:rep}{représentation} à isomorphisme près --- sont déterminées par
le \omterm{def:b3:representations:character}{caractère}. De plus $\langle \chi_V, \chi_V\rangle = \sum_i
m_i^2$~; en particulier $V$ est \omterm{def:b3:representations:rep}{irréductible} ssi
$\langle\chi_V, \chi_V\rangle = 1$.
\end{corollary}

\begin{proof}
$\chi_V = \sum m_i\chi_{V_i}$
(\cref{prop:b3:representations:charbasics})~; prendre les
produits scalaires avec chaque $\chi_{V_i}$ et utiliser
l'orthonormalité. Deux \omterm{def:b3:representations:rep}{représentations} de mêmes \omterm{def:b3:representations:character}{caractères} ont
mêmes multiplicités, donc sont isomorphes.
\end{proof}

\begin{theorem}[La représentation régulière]
\label{thm:b3:representations:regular}
Soient $\chi_1, \dots, \chi_r$ les \omterm{def:b3:representations:character}{caractères} \omterm{def:b3:representations:rep}{irréductibles}
distincts, de \omterm{def:b3:galois:extension}{degrés} $n_i = \chi_i(e)$. La \omterm{ex:b3:representations:examples}{représentation
régulière} se décompose avec multiplicités $m_i = n_i$~; par
suite
\[
\sum_{i=1}^{r} n_i^2 = \abs G,
\qquad
\sum_i n_i\chi_i(g) = 0 \quad (g \neq e).
\]
\end{theorem}

\begin{proof}
Le \omterm{def:b3:representations:character}{caractère} régulier~: $\chi_{\mathrm{reg}}(g) = \#\{h : gh =
h\}$, qui vaut $\abs G$ pour $g = e$ et $0$ sinon. Donc $m_i =
\langle \chi_i,\chi_{\mathrm{reg}}\rangle =
\frac1{\abs G}\overline{\chi_i(e)}\,\abs G = n_i$. Évaluer
$\chi_{\mathrm{reg}} = \sum n_i\chi_i$ en $e$ et en $g \neq e$
donne les deux identités affichées.
\end{proof}

\begin{theorem}\label{thm:b3:representations:numberirr}
Les \omterm{def:b3:representations:character}{caractères} \omterm{def:b3:representations:rep}{irréductibles} forment une \emph{base}
orthonormée de $\mathcal{CF}(G)$~: le nombre de \omterm{def:b3:representations:rep}{représentations
irréductibles} égale le nombre de \omterm{ex:b3:groups:actions}{classes de conjugaison} de $G$.
\end{theorem}

\begin{proof}
Il ne reste que la complétude~: soit $f \in \mathcal{CF}(G)$
orthogonal à tout $\chi_i$~; on montre $f = 0$. Pour une
\omterm{def:b3:representations:rep}{représentation} $(V,\rho)$, poser $T_{f,\rho} =
\frac{1}{\abs G} \sum_g \overline{f(g)}\,\rho(g)$. Elle est
équivariante~: pour $h \in G$,
\[
\rho(h)T_{f,\rho}\rho(h)^{-1} = \frac1{\abs G}\sum_g
\overline{f(g)}\rho(hgh^{-1})
= \frac1{\abs G}\sum_{g'}\overline{f(h^{-1}g'h)}\rho(g') =
T_{f,\rho}
\]
($f$ est une \omterm{def:b3:representations:character}{fonction de classe}). Si $\rho$ est \omterm{def:b3:representations:rep}{irréductible} de
\omterm{def:b3:galois:extension}{degré} $n$, Schur donne $T_{f,\rho} = \lambda\,\mathrm{id}$ avec
\[
\lambda = \frac{\operatorname{tr}T_{f,\rho}}{n}
= \frac{1}{n\abs G}\sum_g \overline{f(g)}\,\chi_\rho(g)
= \frac1n\,\langle f, \chi_\rho\rangle = 0 .
\]
Donc $T_{f,\rho} = 0$ sur toute \omterm{def:b3:representations:rep}{irréductible}, d'où (sommes
directes) sur \emph{toute} \omterm{def:b3:representations:rep}{représentation} --- en particulier
sur la régulière. L'appliquer au vecteur de base $e_e$~: $0 =
T_{f,\mathrm{reg}}e_e = \frac1{\abs
G}\sum_g\overline{f(g)}e_g$, forçant chaque $\overline{f(g)} =
0$. Ainsi la famille orthonormée $(\chi_i)$ engendre
$\mathcal{CF}(G)$, dont la dimension est le nombre de \omterm{ex:b3:groups:actions}{classes
de conjugaison}.
\end{proof}

\begin{corollary}[Orthogonalité des colonnes]
\label{cor:b3:representations:column}
Pour $g, h \in G$~:
\[
\sum_{i=1}^{r}\overline{\chi_i(g)}\,\chi_i(h) =
\begin{cases}
\abs{Z_G(g)} & \text{si } g, h \text{ sont conjugués},\\
0 & \text{sinon.}
\end{cases}
\]
\end{corollary}

\begin{proof}
Soient $g_1, \dots, g_r$ des représentants des classes, $c_j$
les tailles de classes. La matrice $r \times r$ $U_{ij} =
\sqrt{c_j/\abs G}\;\chi_i(g_j)$ a des lignes orthonormées
(\cref{thm:b3:representations:orthogonality} écrite par
classes~: $\sum_j \frac{c_j}{\abs
G}\chi_i(g_j)\overline{\chi_{i'}(g_j)} = \delta_{ii'}$), i.e.\
$UU^* = I$~; une matrice carrée avec $UU^* = I$ a aussi $U^*U =
I$~: les colonnes sont orthonormées, ce qui se déplie en
l'identité affichée ($\abs G/c_j = \abs{Z_G(g_j)}$, la relation
orbite--stabilisateur pour la conjugaison).
\end{proof}

\begin{proposition}[Caractères de dimension un~; relèvement]
\label{prop:b3:representations:onedim}
(a) $G$ est abélien ssi toutes ses \omterm{def:b3:representations:rep}{représentations
irréductibles} sont de \omterm{def:b3:galois:extension}{degré} $1$~; le nombre de \omterm{def:b3:representations:character}{caractères} de
\omterm{def:b3:galois:extension}{degré} $1$ de tout $G$ est $[G : D(G)]$ (ce sont les \omterm{def:b3:representations:character}{caractères}
de l'abélianisé).
(b) Si $N \trianglelefteq G$, les \omterm{def:b3:representations:rep}{représentations irréductibles}
de $G/N$ se relèvent (composer avec $G \to G/N$) en
exactement les \omterm{def:b3:representations:rep}{représentations irréductibles} de $G$ dont le
noyau contient $N$.
\end{proposition}

\begin{proof}
(a) Si $G$ est abélien, chaque classe est un singleton~: $r =
\abs G$, et $\sum n_i^2 = \abs G$ force tous les $n_i = 1$~;
réciproquement si tous les $n_i = 1$, la \omterm{ex:b3:representations:examples}{représentation
régulière} est somme d'unidimensionnelles, donc
$\rho_{\mathrm{reg}}(G)$ est simultanément diagonalisable,
d'où commutative, et $\rho_{\mathrm{reg}}$ est fidèle~: $G$
abélien. Les \omterm{def:b3:representations:rep}{représentations} de \omterm{def:b3:galois:extension}{degré} $1$ sont des morphismes
$G \to \C^\times$ à but abélien~: elles se factorisent par
$G^{\mathrm{ab}} = G/D(G)$ (\cref{exo:b3:groups:9}), et les
\omterm{def:b3:representations:character}{caractères} distincts de l'abélien $G^{\mathrm{ab}}$ sont au
nombre de $\abs{G^{\mathrm{ab}}}$ (autant de classes, tous
\omterm{def:b3:galois:extension}{degrés} $1$). (b) Composer avec la projection préserve
l'irréductibilité (les sous-espaces invariants se
correspondent), et une \omterm{def:b3:representations:rep}{représentation} triviale sur $N$ se
factorise par le quotient (\cref{thm:b3:groups:firstiso}).
\end{proof}

\section{Tables de caractères}

\begin{definition}\label{def:b3:representations:table}
La \emph{table de caractères}\index{table de caractères} de $G$
est la matrice $r \times r$ $\bigl(\chi_i(g_j)\bigr)$~: lignes
indexées par les \omterm{def:b3:representations:character}{caractères} \omterm{def:b3:representations:rep}{irréductibles}, colonnes par les
\omterm{ex:b3:groups:actions}{classes de conjugaison} (avec leurs tailles affichées). Les
lignes sont orthonormées pour le produit pondéré, les colonnes
orthogonales (\cref{cor:b3:representations:column})~: la table
est sévèrement surdéterminée, ce qui la rend calculable.
\end{definition}

\begin{example}[La table de $S_3$]\label{ex:b3:representations:s3}
Classes~: $e$ (taille $1$), transpositions ($3$), $3$-cycles
($2$)~; donc $r = 3$ \omterm{def:b3:representations:rep}{irréductibles}, de \omterm{def:b3:galois:extension}{degrés} $n_i$ avec $\sum
n_i^2 = 6$~: $1, 1, 2$. \omterm{def:b3:galois:extension}{Degré} $1$~: triviale $\mathbf 1$ et
signature $\varepsilon$. La dernière ligne suit de
l'orthogonalité des colonnes (ou de $\chi_{\mathrm{std}} =
\chi_{\mathrm{perm}} - \mathbf 1$)~:
\[
\begin{array}{c|ccc}
S_3 & e & (1\,2)\ [3] & (1\,2\,3)\ [2]\\
\hline
\mathbf 1 & 1 & 1 & 1\\
\varepsilon & 1 & -1 & 1\\
\chi_{\mathrm{std}} & 2 & 0 & -1
\end{array}
\]
Vérification~: $\langle\chi_{\mathrm{std}},\chi_{\mathrm{std}}
\rangle = \frac{1}{6}(4 + 0 + 2) = 1$~: \omterm{def:b3:representations:rep}{irréductible}.
\end{example}

\begin{example}[La table de $S_4$]\label{ex:b3:representations:s4}
Classes~: $e$ [1], transpositions [6], doubles transpositions
[3], $3$-cycles [8], $4$-cycles [6]~: cinq \omterm{def:b3:representations:rep}{irréductibles}, $\sum
n_i^2 = 24$ avec deux de \omterm{def:b3:galois:extension}{degré} $1$ ($\mathbf 1, \varepsilon$~;
$[S_4 : D(S_4)] = [S_4 : A_4] = 2$)~: \omterm{def:b3:galois:extension}{degrés} $1, 1, 2, 3, 3$.
Le \omterm{def:b3:galois:extension}{degré} $2$ se relève de $S_4/V \cong S_3$
(\cref{prop:b3:representations:onedim}(b), $V$ le groupe de
Klein)~; \omterm{def:b3:galois:extension}{degré} $3$~: la \omterm{def:b3:representations:rep}{représentation} standard et sa
\omterm{def:b3:modules:torsion}{torsion} par $\varepsilon$~:
\[
\begin{array}{c|ccccc}
S_4 & e\,[1] & (1\,2)\,[6] & (1\,2)(3\,4)\,[3] &
(1\,2\,3)\,[8] & (1\,2\,3\,4)\,[6]\\
\hline
\mathbf 1 & 1 & 1 & 1 & 1 & 1\\
\varepsilon & 1 & -1 & 1 & 1 & -1\\
\chi_2 & 2 & 0 & 2 & -1 & 0\\
\chi_{\mathrm{std}} & 3 & 1 & -1 & 0 & -1\\
\varepsilon\chi_{\mathrm{std}} & 3 & -1 & -1 & 0 & 1
\end{array}
\]
($\chi_{\mathrm{std}}(g) = \operatorname{fix}(g) - 1$~;
$\chi_2$ évalue la table de $S_3$ sur l'image de chaque classe
mod $V$.) Toutes les vérifications d'orthogonalité de lignes et
colonnes passent --- en faire deux est l'échauffement de
\cref{exo:b3:representations:3}.
\end{example}

\begin{method}\label{met:b3:representations:table}
Pour construire une \omterm{def:b3:representations:table}{table de caractères}~: (1) lister les
\omterm{ex:b3:groups:actions}{classes de conjugaison} et leurs tailles~; (2) compter les
\omterm{def:b3:representations:character}{caractères} de \omterm{def:b3:galois:extension}{degré} $1$ via $G/D(G)$ et les écrire~; (3)
trouver les \omterm{def:b3:galois:extension}{degrés} restants par $\sum n_i^2 = \abs G$ (petite
combinatoire d'entiers)~; (4) obtenir des \omterm{def:b3:representations:rep}{irréductibles}
bon marché~: relever des quotients, soustraire $\mathbf 1$ des
\omterm{def:b3:representations:character}{caractères} de permutation (vérifier $\langle\chi,\chi\rangle =
1$), multiplier des \omterm{def:b3:representations:character}{caractères} connus par des \omterm{def:b3:galois:extension}{degrés} $1$~; (5)
terminer les lignes inconnues par orthogonalité des colonnes
--- chaque colonne est orthogonale aux colonnes déjà
complètes, et la colonne $e$ porte les \omterm{def:b3:galois:extension}{degrés}. Tout vérifier
par un balayage d'orthogonalité complet.
\end{method}

\begin{omfigure}
\begin{tikzpicture}[scale=0.52]
  \draw[thick] (0,0) rectangle (6,6);
  \fill[omDef!25] (0,5) rectangle (1,6);
  \fill[omProp!25] (1,4) rectangle (2,5);
  \fill[omThm!20] (2,0) rectangle (6,4);
  \draw (0,5) rectangle (1,6);
  \draw (1,4) rectangle (2,5);
  \draw (2,0) rectangle (6,4);
  \node[font=\small] at (0.5,5.5) {$1$};
  \node[font=\small] at (1.5,4.5) {$1$};
  \node[font=\small] at (4,2) {$M_2(\C)$};
\end{tikzpicture}

{\small La \omterm{ex:b3:representations:examples}{représentation régulière} de $S_3$, diagonalisée par
blocs~: $\C[S_3] \cong \C \times \C \times M_2(\C)$, dimensions
$1 + 1 + 4 = 6 = \abs{S_3}$. En général $\C[G] \cong \prod_i
M_{n_i}(\C)$~: l'identité $\sum n_i^2 = \abs G$ est un énoncé
sur des blocs matriciels.}
\end{omfigure}

\section{Exercices}

\begin{exercise}[$\star$]\label{exo:b3:representations:1}
(a) Montrer que les \omterm{def:b3:representations:character}{caractères} \omterm{def:b3:representations:rep}{irréductibles} de $\Z/n\Z$ sont
les $\chi_k(\bar m) = \eu^{2\iu\pi km/n}$, $k = 0, \dots, n-1$,
et écrire la \omterm{def:b3:representations:table}{table de caractères} de $\Z/4\Z$.
(b) Vérifier les deux relations d'orthogonalité dessus --- et
reconnaître la matrice~: où ce livre l'a-t-il déjà vue~?
\end{exercise}

\begin{exercise}[$\star$]\label{exo:b3:representations:2}
Soit $G$ agissant sur un ensemble fini $X$ et $\chi$ le
\omterm{def:b3:representations:character}{caractère} de la \omterm{def:b3:representations:rep}{représentation} de permutation $\C^X$.
(a) Montrer $\chi(g) = \abs{\operatorname{Fix}_X(g)}$ et
$\langle \mathbf 1, \chi\rangle = \#\{\text{orbites}\}$ --- le
lemme de dénombrement de Burnside (\cref{exo:b3:groups:5}) est
un calcul de \omterm{def:b3:representations:character}{caractères}.
(b) Supposer l'\omterm{def:b3:groups:action}{action} transitive, donc $\chi = \mathbf 1 +
\psi$. Montrer que $\psi$ est \omterm{def:b3:representations:rep}{irréductible} ssi l'\omterm{def:b3:groups:action}{action} est
\emph{$2$-transitive} (transitive sur les couples ordonnés de
points distincts). \emph{(Calculer $\langle\chi,\chi\rangle$
comme le nombre d'\omterm{def:b3:groups:action}{orbites} sur $X \times X$.)}
(c) Conclure que la \omterm{def:b3:representations:rep}{représentation} standard de $S_n$ ($n \geq
2$) est \omterm{def:b3:representations:rep}{irréductible}.
\end{exercise}

\begin{exercise}[$\star$]\label{exo:b3:representations:3}
Reconstruire la table de $S_3$ depuis zéro en suivant
la~\cref{met:b3:representations:table}, puis vérifier deux
relations d'orthogonalité de lignes et deux de colonnes dans la
table de $S_4$ de \cref{ex:b3:representations:s4}. Décomposer
le \omterm{def:b3:representations:character}{caractère} de permutation de $S_4$ agissant sur $\{1,2,3,4\}$
et le \omterm{def:b3:representations:character}{caractère} $\chi_{\mathrm{std}}^2$ (carré pointwise) en
\omterm{def:b3:representations:rep}{irréductibles}.
\end{exercise}

\begin{exercise}[$\star\star$]\label{exo:b3:representations:4}
Calculer les \omterm{def:b3:representations:table}{tables de caractères} de $D_4$ et de $Q_8$.
Conclure que deux groupes non isomorphes peuvent avoir des
\omterm{def:b3:representations:table}{tables de caractères} identiques --- quelles données
groupes-théoriques la table capture-t-elle néanmoins dans cette
paire (ordres des \omterm{ex:b3:groups:actions}{centres}, abélianisés, nombre d'involutions)~?
Lesquelles \emph{échoue}-t-elle à capturer~?
\end{exercise}

\begin{exercise}[$\star\star$]\label{exo:b3:representations:5}
\omterm{def:b3:representations:table}{Table de caractères} de $A_4$~: classes $e$ [1], doubles
transpositions [3], et \emph{deux} classes de $3$-cycles [4],
[4].
(a) Expliquer le scindage des $3$-cycles (comparer les
\omterm{ex:b3:groups:actions}{centralisateurs} dans $S_4$ et $A_4$, comme dans
\cref{exo:b3:groups:11}).
(b) Trouver les trois \omterm{def:b3:representations:character}{caractères} de \omterm{def:b3:galois:extension}{degré} $1$ (via $A_4/V \cong
\Z/3\Z$) et le \omterm{def:b3:representations:character}{caractère} de \omterm{def:b3:galois:extension}{degré} $3$ (restreindre
$\chi_{\mathrm{std}}$ de $S_4$), et assembler la table.
(c) Lire les sous-groupes normaux de $A_4$ sur la table (noyaux
$\{g : \chi(g) = \chi(e)\}$ et leurs intersections).
\end{exercise}

\begin{exercise}[$\star\star$]\label{exo:b3:representations:6}
(a) Montrer que $\ker\chi = \{g : \chi(g) = \chi(e)\}$ est le
noyau de la \omterm{def:b3:representations:rep}{représentation} sous-jacente
(\cref{prop:b3:representations:charbasics}, cas d'égalité).
(b) Montrer que tout \omterm{def:b3:groups:normal}{sous-groupe normal} de $G$ est une
intersection de noyaux de \omterm{def:b3:representations:character}{caractères} \omterm{def:b3:representations:rep}{irréductibles}.
\emph{(Représenter $G/N$ fidèlement~: sa \omterm{ex:b3:representations:examples}{représentation
régulière}.)}
(c) En déduire~: $G$ est \omterm{def:b3:groups:simple}{simple} ssi $\ker\chi_i = \{e\}$ pour
tout \omterm{def:b3:representations:rep}{irréductible} non trivial $\chi_i$ --- la simplicité se lit
sur la \omterm{def:b3:representations:table}{table de caractères}.
\end{exercise}

\begin{exercise}[$\star\star$]\label{exo:b3:representations:7}
\omterm{def:b3:galois:extension}{Degrés} pour $|G| = 8$~: montrer qu'un groupe non abélien d'ordre
$8$ a le motif de \omterm{def:b3:galois:extension}{degrés} $(1,1,1,1,2)$, et que sa \omterm{def:b3:representations:rep}{représentation}
de \omterm{def:b3:galois:extension}{degré} $2$ est fidèle. Plus généralement montrer qu'un groupe
non abélien d'ordre $p^3$ a le motif $(1^{\,p^2}, p, \dots, p)$
avec $p^2$ uns et $p - 1$ \omterm{def:b3:representations:character}{caractères} de \omterm{def:b3:galois:extension}{degré} $p$.
\emph{(Utiliser $[G : D(G)]$ et $\sum n_i^2 = |G|$~; ici $D(G)
= Z(G)$ a ordre $p$.)}
\end{exercise}

\begin{exercise}[$\star\star\star$]\label{exo:b3:representations:8}
Pour des groupes finis $G, H$~: montrer que les \omterm{def:b3:representations:character}{fonctions de
classe} $\chi(g)\psi(h)$ sur $G \times H$, pour $\chi, \psi$
\omterm{def:b3:representations:character}{caractères} \omterm{def:b3:representations:rep}{irréductibles} de $G, H$, sont exactement les
\omterm{def:b3:representations:character}{caractères} \omterm{def:b3:representations:rep}{irréductibles} de $G \times H$.
\emph{(L'orthonormalité est un calcul direct~; la complétude
en comptant les classes.)} En déduire la \omterm{def:b3:representations:table}{table de caractères} de
$\Z/2\Z \times \Z/2\Z$ et redériver
la~\cref{prop:b3:representations:onedim} pour les groupes
abéliens finis via le théorème de structure.
\end{exercise}

\begin{exercise}[$\star\star\star$]\label{exo:b3:representations:9}
La \omterm{def:b3:representations:table}{table de caractères} de $A_5$ (classes de tailles $1, 15, 20,
12, 12$ de \cref{exo:b3:groups:11})~:
(a) Montrer que les \omterm{def:b3:galois:extension}{degrés} sont $1, 3, 3, 4, 5$ \emph{(la seule
solution de $\sum n_i^2 = 60$ avec $n_1 = 1$ et, en utilisant
\cref{exo:b3:representations:6}(c) avec la simplicité, aucun
autre $n_i = 1$)}.
(b) Construire le \omterm{def:b3:representations:character}{caractère} de \omterm{def:b3:galois:extension}{degré} $4$ (\omterm{def:b3:groups:action}{action} de permutation
sur $5$ points) et le \omterm{def:b3:representations:character}{caractère} de \omterm{def:b3:galois:extension}{degré} $5$ (l'\omterm{def:b3:groups:action}{action} sur les
six $5$-sous-groupes de Sylow donne le \omterm{def:b3:galois:extension}{degré} $6 = 1 + 5$~;
vérifier l'irréductibilité), et compléter les deux lignes de
\omterm{def:b3:galois:extension}{degré} $3$ par orthogonalité des colonnes~: des entrées en
nombre d'or $\frac{1\pm\sqrt5}2$ apparaissent sur les classes
de $5$-cycles.
(c) Vérifier sur la table finie que $A_5$ est \omterm{def:b3:groups:simple}{simple}
(\cref{exo:b3:representations:6}(c)).
\end{exercise}

\begin{exercise}[$\star\star$]\label{exo:b3:representations:10}
Soit $\rho$ une \omterm{def:b3:representations:rep}{représentation irréductible} de \omterm{def:b3:galois:extension}{degré} $n$ et $z
\in Z(G)$. Montrer $\rho(z) = \lambda_z\,\mathrm{id}$ avec
$\lambda \colon Z(G) \to \C^\times$ un morphisme (le
\emph{\omterm{def:b3:representations:character}{caractère} central}), et en déduire $|\chi(z)| = n$ pour
$z$ central. Application~: si $G$ a une \omterm{def:b3:representations:rep}{représentation
irréductible} fidèle, alors $Z(G)$ est cyclique.
\end{exercise}

\begin{exercise}[$\star\star$]\label{exo:b3:representations:11}
(Projections isotypiques) Soit $(V, \rho)$ une \omterm{def:b3:representations:rep}{représentation} de
$G$ et $\chi_i$ un \omterm{def:b3:representations:character}{caractère} \omterm{def:b3:representations:rep}{irréductible} de \omterm{def:b3:galois:extension}{degré} $n_i$.
Définir
\[
p_i = \frac{n_i}{\abs G}\sum_{g\in G}
\overline{\chi_i(g)}\,\rho(g) \;\in\; \mathcal L(V).
\]
(a) Montrer que $p_i$ est $G$-équivariant, et calculer sa
restriction à une \omterm{def:b3:representations:rep}{sous-représentation} \omterm{def:b3:representations:rep}{irréductible} $W
\subseteq V$ de \omterm{def:b3:representations:character}{caractère} $\chi_j$~: c'est $\delta_{ij}\,
\mathrm{id}_W$ \emph{(Schur~; prendre les traces pour
identifier le scalaire)}.
(b) En déduire que $p_i$ est une projection sur la somme $V_i$
de toutes les \omterm{def:b3:representations:rep}{sous-représentations} \omterm{def:b3:representations:rep}{irréductibles} de \omterm{def:b3:representations:character}{caractère}
$\chi_i$ (la \emph{composante isotypique}), que $\sum_ip_i =
\mathrm{id}_V$, et que la décomposition $V = \bigoplus_iV_i$
est canonique --- contrairement au scindage plus fin de chaque
$V_i$ en \omterm{def:b3:representations:rep}{irréductibles}.
(c) Pour la \omterm{ex:b3:representations:examples}{représentation régulière} de $S_3$ et le \omterm{def:b3:representations:character}{caractère}
signature $\varepsilon$, écrire $p_\varepsilon$ explicitement
comme élément de l'algèbre du groupe et vérifier
$p_\varepsilon^2 = p_\varepsilon$ à la main.
\end{exercise}

\begin{exercise}[$\star\star$]\label{exo:b3:representations:12}
(Lire une table) La \omterm{def:b3:representations:table}{table de caractères} d'un certain groupe $G$
d'ordre $24$ est partiellement connue~: elle a $5$ classes, de
tailles $1, 6, 8, 6, 3$, et des \omterm{def:b3:galois:extension}{degrés} $1, 1, 2, 3, 3$.
(a) Retrouver la table complète~: les deux \omterm{def:b3:representations:character}{caractères} linéaires
(un trivial~; l'autre prend la valeur $-1$ exactement sur les
classes de tailles $6$ et $6$), puis le \omterm{def:b3:representations:character}{caractère} de \omterm{def:b3:galois:extension}{degré} $2$
via l'orthogonalité des colonnes avec la colonne de l'identité,
puis les deux \omterm{def:b3:representations:character}{caractères} de \omterm{def:b3:galois:extension}{degré} $3$ de même (l'un est
$\chi_2\chi_4$).
(b) Identifier $G$ ($\cong S_4$~: comparer les classes aux
types de cycles), et extraire de la table les sous-groupes
normaux via \cref{exo:b3:representations:6}~: noyaux de $\chi_2$
(indice $2$~: $A_4$) et de $\chi_3$ (le Klein $V_4$), et rien
d'autre que $\{e\}, G$.
(c) Expliquer comment la table montre $G/V_4 \cong S_3$
\emph{(quels \omterm{def:b3:representations:character}{caractères} se factorisent à travers le
quotient~?)}.
\end{exercise}

\section{Problème~: le théorème \texorpdfstring{$p^aq^b$}{paqb}
de Burnside}

\begin{problem}[{Problème du week-end --- résolubilité des
groupes d'ordre $p^aq^b$}]\label{pb:b3:representations:1}
Burnside montra en 1904 que tout groupe dont l'ordre a au plus
deux facteurs premiers est \omterm{def:b3:groups:derived}{résoluble} --- un énoncé sur les
groupes abstraits dont les seules démonstrations connues pendant
un demi-siècle passaient par la théorie des \omterm{def:b3:representations:character}{caractères}. Ce
problème construit la démonstration en entier, en assemblant
le~\cref{ch:b3:groups} (résolubilité),
le~\cref{ch:b3:modules} (\omterm{def:b3:modules:free}{modules de type fini} sur $\Z$) et ce
chapitre. Tout au long, $\chi_1, \dots, \chi_r$ sont les
\omterm{def:b3:representations:character}{caractères} \omterm{def:b3:representations:rep}{irréductibles} de $G$, $n_i = \chi_i(e)$.

\medskip
\noindent\textbf{Partie I --- Entiers \omterm{def:b3:galois:algebraic}{algébriques}.} Un
\emph{entier \omterm{def:b3:galois:algebraic}{algébrique}} est une racine d'un polynôme
\emph{unitaire} de $\Z[X]$.
\begin{enumerate}
  \item Montrer que $\alpha$ est un entier \omterm{def:b3:galois:algebraic}{algébrique} ssi
        l'anneau $\Z[\alpha]$ est un $\Z$-module \omterm{def:b3:modules:free}{de type fini}.
  \item En déduire que les entiers \omterm{def:b3:galois:algebraic}{algébriques} forment un
        sous-anneau de $\C$. \emph{(Si $\Z[\alpha], \Z[\beta]$
        sont \omterm{def:b3:modules:free}{de type fini}, $\Z[\alpha, \beta]$ l'est aussi, et
        les \omterm{def:b3:modules:module}{sous-modules} des $\Z$-modules \omterm{def:b3:modules:free}{de type fini} sont \omterm{def:b3:modules:free}{de
        type fini}, par le~\cref{thm:b3:modules:submodule}.)}
  \item Montrer qu'un entier \omterm{def:b3:galois:algebraic}{algébrique} \emph{rationnel} est un
        entier. \emph{(Théorème des racines rationnelles.)}
  \item Montrer que toute valeur de \omterm{def:b3:representations:character}{caractère} $\chi(g)$ est un
        entier \omterm{def:b3:galois:algebraic}{algébrique}.
\end{enumerate}

\medskip
\noindent\textbf{Partie II --- Les relations de sommes de
classes.} Fixer une \omterm{def:b3:representations:rep}{irréductible} $(V, \rho)$ de \omterm{def:b3:galois:extension}{degré} $n$ et de
\omterm{def:b3:representations:character}{caractère} $\chi$. Pour une \omterm{ex:b3:groups:actions}{classe de conjugaison} $C$, soit $S_C
= \sum_{g \in C}\rho(g) \in \mathcal L(V)$.
\begin{enumerate}[resume]
  \item Montrer que $S_C$ est équivariant, donc $S_C =
        \omega_C\,\mathrm{id}$ avec
        \[
        \omega_C = \frac{\abs C\,\chi(g_C)}{n}
        \qquad (g_C \in C \text{ un représentant quelconque}).
        \]
  \item Montrer que $S_CS_{C'} = \sum_{C''}
        a_{CC'C''}\,S_{C''}$ où $a_{CC'C''} \in \N$ compte,
        pour un $z \in C''$ fixé, les paires $(x, y) \in C
        \times C'$ avec $xy = z$. En déduire que les $\omega_C$
        vérifient $\omega_C\,\omega_{C'} = \sum_{C''}
        a_{CC'C''}\,\omega_{C''}$.
  \item Conclure que chaque $\omega_C$ est un entier
        \omterm{def:b3:galois:algebraic}{algébrique}.
  \item En déduire la \emph{divisibilité de \omterm{thm:b3:galois:finitefields}{Frobenius}}~: $n_i$
        divise $\abs G$ pour tout \omterm{def:b3:galois:extension}{degré} \omterm{def:b3:representations:rep}{irréductible}.
        \emph{(Calculer $\frac{\abs G}{n}$ comme entier
        \omterm{def:b3:galois:algebraic}{algébrique} rationnel.)}
\end{enumerate}

\medskip
\noindent\textbf{Partie III --- Le critère de simplicité de
Burnside.}
\begin{enumerate}[resume]
  \item Soit $\chi$ \omterm{def:b3:representations:rep}{irréductible} de \omterm{def:b3:galois:extension}{degré} $n$ et $C$ une classe
        avec $\gcd(\abs C, n) = 1$. En utilisant Bézout et les
        questions 4--7, montrer que $\frac{\chi(g_C)}{n}$ est
        un entier \omterm{def:b3:galois:algebraic}{algébrique}.
  \item Supposer de plus $0 < \abs{\chi(g_C)} < n$. Montrer que
        c'est impossible~: l'entier \omterm{def:b3:galois:algebraic}{algébrique} $\alpha =
        \chi(g_C)/n$ a tous ses \emph{conjugués} de \omterm{def:b3:modules:module}{module}
        $\leq 1$, donc le produit des conjugués $N$ est un
        entier \omterm{def:b3:galois:algebraic}{algébrique} rationnel avec $0 < \abs N < 1$.
        Conclure~: soit $\chi(g_C) = 0$, soit $\rho(g_C)$ est
        scalaire.
  \item (Critère de Burnside) Soit $C \neq \{e\}$ une \omterm{ex:b3:groups:actions}{classe de
        conjugaison} de taille \emph{puissance de premier} $p^k
        > 1$, et supposer $G$ \omterm{def:b3:groups:simple}{simple} non abélien. L'orthogonalité
        des colonnes sur la colonne de $C$ contre celle de $e$
        donne $1 + \sum_{i \geq 2} n_i\chi_i(g_C) = 0$. Montrer
        qu'un certain $\chi_i$ non trivial avec $p \nmid n_i$ a
        $\chi_i(g_C) \neq 0$~; par la question 10, $\rho_i(g_C)$
        est scalaire~; dériver une contradiction avec la
        simplicité. Conclure~: \emph{aucun \omterm{def:b3:groups:simple}{groupe simple} non
        abélien n'a de \omterm{ex:b3:groups:actions}{classe de conjugaison} de taille puissance
        de premier $> 1$.}
\end{enumerate}

\medskip
\noindent\textbf{Partie IV --- Le théorème.}
\begin{enumerate}[resume]
  \item Soit $\abs G = p^aq^b$ avec $a + b \geq 1$. Si $G$ est
        \omterm{def:b3:groups:simple}{simple}, montrer qu'il est abélien~: choisir $z \neq e$
        dans le \omterm{ex:b3:groups:actions}{centre} d'un $q$-sous-groupe de Sylow
        (\cref{thm:b3:groups:pfixed}) et considérer la taille
        de sa \omterm{ex:b3:groups:actions}{classe de conjugaison} $[G : Z_G(z)]$, une
        puissance de $p$~; appliquer la question 11.
  \item Conclure par récurrence sur $\abs G$~: \emph{tout
        groupe d'ordre $p^aq^b$ est \omterm{def:b3:groups:derived}{résoluble}} (Burnside).
        Pourquoi l'argument casse-t-il pour trois premiers ---
        et le doit-il, étant donné $\abs{A_5} = 2^2\cdot3\cdot
        5$~?
\end{enumerate}

\medskip
\noindent\textbf{Partie V --- La \omterm{def:b3:representations:table}{table de caractères} de $A_5$.} Le
plus petit groupe que le théorème de Burnside ne peut toucher
mérite son portrait complet~; tout ce qui suit n'utilise que ce
chapitre et \cref{exo:b3:groups:11}.
\begin{enumerate}[resume]
  \item Rappeler de \cref{exo:b3:groups:11} les cinq \omterm{ex:b3:groups:actions}{classes de
        conjugaison} de $A_5$~: $\{e\}$, les $15$ doubles
        transpositions, les $20$ $3$-cycles, et deux classes de
        $12$ $5$-cycles chacune, représentées par $c =
        (1\,2\,3\,4\,5)$ et $c^2$. Expliquer pourquoi les
        $5$-cycles se scindent en deux classes de $A_5$ alors
        qu'ils forment une seule classe de $S_5$.
  \item Montrer que les \omterm{def:b3:galois:extension}{degrés} \omterm{def:b3:representations:rep}{irréductibles} de $A_5$ sont
        exactement $1, 3, 3, 4, 5$~: utiliser $\sum_in_i^2 = 60$
        avec $r = 5$ classes, et le fait que $A_5$ est \omterm{prop:b3:galois:perfect}{parfait}
        ($D(A_5) = A_5$), donc le \omterm{def:b3:representations:character}{caractère} trivial est son
        seul linéaire~; puis éliminer tout autre multiensemble
        \emph{(écrire $59$ comme somme de quatre carrés d'entiers
        $\geq 2$~: vérifier qu'il n'y a qu'une façon avec des
        \omterm{def:b3:galois:extension}{degrés} plausibles)}.
  \item Soit $\pi$ le \omterm{def:b3:representations:character}{caractère} de permutation de $A_5$ sur
        $\{1, \dots, 5\}$~: $\pi(g) = \#\operatorname{Fix}(g)$,
        avec valeurs $5, 1, 2, 0, 0$ sur les cinq classes.
        Calculer $\langle\pi, \mathbf 1\rangle$ et $\langle\pi,
        \pi\rangle$, et en déduire que $\chi_4 = \pi - \mathbf 1$
        est \omterm{def:b3:representations:rep}{irréductible} de \omterm{def:b3:galois:extension}{degré} $4$, avec valeurs $4, 0, 1,
        -1, -1$.
  \item Même jeu sur les $10$ paires non ordonnées $\{i, j\}$~:
        les comptages de points fixes sont $10, 2, 1, 0, 0$.
        Calculer $\langle\pi_{10}, \pi_{10}\rangle$,
        $\langle\pi_{10}, \mathbf 1\rangle$ et
        $\langle\pi_{10}, \chi_4\rangle$, en déduire la
        décomposition $\pi_{10} = \mathbf 1 + \chi_4 +
        \chi_5$, et obtenir l'\omterm{def:b3:representations:rep}{irréductible} $\chi_5$ de \omterm{def:b3:galois:extension}{degré} $5$
        avec valeurs $5, 1, -1, 0, 0$.
  \item Les deux \omterm{def:b3:representations:rep}{irréductibles} restants $\chi_2, \chi_3$ ont
        \omterm{def:b3:galois:extension}{degré} $3$. L'orthogonalité des colonnes (chaque colonne
        non identité contre celle de l'identité, et chaque
        colonne avec elle-même) détermine leurs valeurs hors des
        $5$-cycles~: montrer $\chi_2(g) = \chi_3(g) = -1$ sur
        les doubles transpositions et $0$ sur les $3$-cycles.
  \item Sur les classes de $5$-cycles, poser $x = \chi_2(c)$ et
        $y = \chi_2(c^2)$~; par symétrie on peut prendre
        $\chi_3(c) = y$, $\chi_3(c^2) = x$. De la colonne de $c$
        appairée avec celle de l'identité et avec celle de
        $c^2$, dériver $x + y = 1$ et $xy = -1$ (et vérifier la
        valeur $x^2 + y^2 = 3$ donnée par la colonne de $c$
        avec elle-même), d'où
        \[
        \{x, y\} = \Bigl\{\frac{1 + \sqrt5}2,\ \frac{1 -
        \sqrt5}2\Bigr\} :
        \]
        le nombre d'or et son conjugué. Assembler la \omterm{def:b3:representations:table}{table de
        caractères} complète de $A_5$.
  \item Lancer les vérifications~: la norme de ligne de $\chi_2$
        vaut $1$ (utiliser $\varphi^2 + \bar\varphi^2 = 3$),
        $\langle\chi_2, \chi_3\rangle = 0$, et la divisibilité
        de \omterm{thm:b3:galois:finitefields}{Frobenius} (question 8) pour les cinq \omterm{def:b3:galois:extension}{degrés}. Où dans
        la table voit-on une différence avec $S_5$, dont toutes
        les valeurs de \omterm{def:b3:representations:character}{caractères} sont des entiers
        rationnels~?
  \item Déduire de la table seule que $A_5$ est \omterm{def:b3:groups:simple}{simple}~: un
        \omterm{def:b3:groups:normal}{sous-groupe normal} est une réunion de \omterm{ex:b3:groups:actions}{classes de
        conjugaison} contenant $e$ dont le cardinal divise $60$
        --- vérifier qu'aucune sous-somme propre de $1 + 15 +
        20 + 12 + 12$ contenant le terme $1$ ne divise $60$.
        Recouper avec le critère de la question 11~: vérifier
        qu'aucune classe de $A_5$ n'a taille puissance de
        premier $> 1$.
  \item (Coda icosaédrale) $A_5$ est le groupe des rotations de
        l'icosaèdre, et les \omterm{def:b3:representations:rep}{représentations} de \omterm{def:b3:galois:extension}{degré} $3$ sont
        les deux \omterm{def:b3:groups:action}{actions} géométriques sur $\R^3$. Vérifier
        l'identité de trace~: une rotation d'angle $\theta$ a
        pour trace $1 + 2\cos\theta$, et $1 +
        2\cos\frac{2\pi}5 = \frac{1+\sqrt5}2 = \varphi$.
        Expliquer sans calcul pourquoi l'autre \omterm{def:b3:representations:character}{caractère} de
        \omterm{def:b3:galois:extension}{degré} $3$ doit porter la valeur conjuguée~: le \omterm{def:b3:galois:galois}{groupe
        de Galois} de $\Q(\sqrt5)/\Q$ agit sur toute la \omterm{def:b3:representations:table}{table de
        caractères} (entrée par entrée), en permutant les
        \omterm{def:b3:representations:character}{caractères} \omterm{def:b3:representations:rep}{irréductibles}.
\end{enumerate}

\medskip
\noindent\textbf{Partie VI --- Compléments~: une borne centrale
et le carré tensoriel.}
\begin{enumerate}[resume]
  \item (Plus fin qu'une divisibilité) Soit $\chi$ \omterm{def:b3:representations:rep}{irréductible}
        de \omterm{def:b3:galois:extension}{degré} $n$. Montrer que $\abs{\chi(z)} = n$ pour tout
        $z \in Z(G)$ \emph{(lemme de Schur~: $\rho(z)$ est
        scalaire, d'ordre fini)}, et déduire de $\langle\chi,
        \chi\rangle = 1$ la borne
        \[
        n^2 \leq [G : Z(G)] .
        \]
        Montrer que les groupes non abéliens d'ordre $8$ ont
        des \omterm{def:b3:galois:extension}{degrés} \omterm{def:b3:representations:rep}{irréductibles} $1, 1, 1, 1, 2$ \emph{(cinq
        \omterm{ex:b3:groups:actions}{classes de conjugaison}~; écrire $8$ comme somme de
        cinq carrés)} et atteignent l'égalité $4 = [G :
        Z(G)]$~; vérifier la borne sur $A_5$, dont le \omterm{ex:b3:groups:actions}{centre}
        est trivial.
  \item (Carré tensoriel de $\chi_2$) Pour $g$ d'ordre fini,
        $\rho(g)$ est diagonalisable à valeurs propres racines
        de l'unité~; en déduire les formules de \omterm{def:b3:representations:character}{caractères}
        \[
        \chi_{\operatorname{Sym}^2 V}(g)
        = \frac{\chi(g)^2 + \chi(g^2)}2,
        \qquad
        \chi_{\Lambda^2 V}(g)
        = \frac{\chi(g)^2 - \chi(g^2)}2 .
        \]
        Les appliquer à $\chi_2$ de $A_5$ \emph{(noter que $g^2$
        parcourt la classe de $c^2$ quand $g$ parcourt celle de
        $c$, et réciproquement)}~: montrer $\Lambda^2\chi_2 =
        \chi_2$ et $\operatorname{Sym}^2\chi_2 = \mathbf 1 +
        \chi_5$, d'où
        \[
        \chi_2\otimes\chi_2 = \mathbf 1 + \chi_2 + \chi_5 .
        \]
        Interpréter $\Lambda^2\chi_2 = \chi_2$ géométriquement
        via le produit vectoriel sur $\R^3$.
  \item (Audit final de la table) Vérifier numériquement~:
        l'orthogonalité des colonnes entre les deux colonnes de
        $5$-cycles ($\varphi\bar\varphi = -1$), la valeur $5 =
        \abs{Z_{A_5}(c)}$ pour la colonne de $c$ contre
        elle-même, et l'annulation du \omterm{def:b3:representations:character}{caractère} régulier
        $\sum_i n_i\chi_i(g) = 0$ sur chacune des quatre
        colonnes non identité de la table.
\end{enumerate}
\end{problem}
